% Problem record op_4154bf4cbe0d288e. This TeX file is derived from
% database/problems_json/op_4154bf4cbe0d288e.json by scripts/sync-tex.mjs; edit the JSON record
% and rerun the script rather than editing this file.
\section{Deterministic quadratic-size vertex-minor-universal graphs}
\label{sec:op_4154bf4cbe0d288e}

\paragraph{Problem.}
Does there exist an absolute constant $C>0$ and a deterministic algorithm that, for each integer $k\geq2$ supplied in unary, runs in time $k^{O(1)}$ and outputs a $k$-vertex-minor-universal graph with at most $Ck^2$ vertices? The output is a finite simple labelled graph $G_k=(V_k,E_k)$. Write $H\leq_{\mathrm{vm}}G_k$ when $H$ is obtainable by local complementations and vertex deletions that preserve the labels of surviving vertices; local complementation toggles edges between distinct neighbors of a vertex. Universality requires every graph on every prescribed $k$-element subset of $V_k$ to be obtainable, as in Eq.~\eqref{eq:gvu-universal}:
\begin{equation}
\begin{gathered}
k\leq|V_k|\leq Ck^2,\\
\forall S\subseteq V_k\text{ with }|S|=k,\quad
\forall E_H\subseteq\{\{u,v\}:u,v\in S,\ u\neq v\},\\
(S,E_H)\leq_{\mathrm{vm}}G_k.
\end{gathered}
\label{eq:gvu-universal}
\end{equation}

\subsection*{Status}

Unsolved

\subsection*{Source}

This is a precise algorithmic formulation of the first open question in Section 5, p.~36:15 of Cautrès et al., which asks for deterministic cubic or quadratic constructions of vertex-minor-universal graphs \sourcecite{ref:gvu-universal}{CCM+24}. The unary input and polynomial-time requirement make explicit the efficient-construction target; the present question selects the quadratic-size case.

\subsection*{Progress}

\begin{itemize}
  \item For an $n$-vertex $k$-vertex-minor-universal graph, counting local-Clifford classes (Proposition 5 and its proof) obtainable by Pauli measurements gives
  \begin{equation}
3^{n-k}\geq2^{(k^2-5k)/2-1},
  \qquad
  n\geq k+\frac{k^2-5k-2}{2\log_2 3}
  =\Omega(k^2).
\label{eq:gvu-progress-2}
\end{equation}
  Equation~\eqref{eq:gvu-progress-2} shows that quadratic order is necessary, irrespective of construction time. \sourcecite{ref:gvu-small-pairable}{CMP23}

  \item Theorem 8 of Cautrès et al. establishes that, for every real $\alpha>2$ and all sufficiently large $k$, there exists a $k$-vertex-minor-universal graph satisfying
  \begin{equation}
|V_k|\leq\alpha k^2.
\label{eq:gvu-progress-3}
\end{equation}
  By Eq.~\eqref{eq:gvu-progress-3}, the optimal asymptotic order is already known existentially. \sourcecite{ref:gvu-universal}{CCM+24}

  \item For a prime power $q$, let $B_q$ be the bipartite incidence graph between points and lines of the projective plane over the field with $q$ elements; the explicit construction in Theorem 19 satisfies
  \begin{equation}
|V(B_q)|=2(q^2+q+1),
  \qquad
  7k^2-16\leq4q
  \quad\Longrightarrow\quad
  B_q\text{ is }k\text{-vertex-minor universal}.
\label{eq:gvu-progress-4}
\end{equation}
  In Eq.~\eqref{eq:gvu-progress-4}, choosing $q=\Theta(k^2)$ gives a deterministic polynomial-time construction with $O(k^4)$ vertices; the authors explicitly ask for cubic or quadratic deterministic constructions in their conclusion. \sourcecite{ref:gvu-universal}{CCM+24}

  \item Ascoli and coauthors improve (Theorem 1.1) the random-graph guarantee: writing $\mathbb G(n,p)$ for the distribution in which each possible edge appears independently with probability $p$, and defining $c_0:=1/(2\log_2(4/3))$, for every fixed $\eta>0$ one has, asymptotically as $k\to\infty$,
  \begin{equation}
n\geq(1+\eta)c_0k^2
  \quad\Longrightarrow\quad
  \Pr_{G\sim\mathbb G(n,1/2)}
  [G\text{ is }k\text{-vertex-minor universal}]
  \geq1-2^{-(1+o(1))\eta k^2/2}.
\label{eq:gvu-progress-5}
\end{equation}
  Equation~\eqref{eq:gvu-progress-5} gives efficient random generation with high success probability, not a guaranteed deterministic construction. \sourcecite{ref:gvu-random}{AFF+26}

  \item Chao and Xu's May 2026 revision (version 3, Theorems 1.3 and 1.4, and Appendix A) resolves the associated random-graph universality conjecture across its probability regimes; in particular, for sufficiently large $n$, with $s:=\min(p,1-p)$,
  \begin{equation}
\begin{gathered}
s\geq\frac{100\log_2 n}{\sqrt n},
  \\

  k\leq\frac{s\sqrt n}{100}
  \\
\Longrightarrow\quad
  \Pr_{G\sim\mathbb G(n,p)}
  [G\text{ is }k\text{-vertex-minor universal}]
  \geq1-2^{-s^2n/100}.
\end{gathered}
\label{eq:gvu-progress-6}
\end{equation}
  Beyond the regime in Eq.~\eqref{eq:gvu-progress-6}, their additional sparse and dense regime result completes the probabilistic conjecture, but does not supply the deterministic quadratic-size family asked for here. \sourcecite{ref:gvu-random-regimes}{CX26}
\end{itemize}

\subsection*{References}

\begin{enumerate}[label={},leftmargin=4.8em,labelsep=0.5em]
  \footnotesize
  \item[\textup{[CMP23]}]\label{ref:gvu-small-pairable}
  N. Claudet, M. Mhalla, and S. Perdrix, ``Small $k$-Pairable States,'' arXiv preprint (2023). \href{https://doi.org/10.48550/arXiv.2309.09956}{doi:10.48550/arXiv.2309.09956}; \href{https://arxiv.org/abs/2309.09956}{arXiv:2309.09956}.

  \item[\textup{[CCM+24]}]\label{ref:gvu-universal}
  M. Cautrès, N. Claudet, M. Mhalla, S. Perdrix, V. Savin, and S. Thomassé, ``Vertex-Minor Universal Graphs for Generating Entangled Quantum Subsystems,'' in \emph{51st International Colloquium on Automata, Languages, and Programming (ICALP 2024)}, 36:1–36:18 (2024). \href{https://doi.org/10.4230/LIPIcs.ICALP.2024.36}{doi:10.4230/LIPIcs.ICALP.2024.36}; \href{https://arxiv.org/abs/2402.06260}{arXiv:2402.06260}.

  \item[\textup{[AFF+26]}]\label{ref:gvu-random}
  R. Ascoli, B. Frederickson, S. Frederickson, C. McFarland, and L. Post, ``Almost All Graphs Are Vertex-Minor Universal,'' arXiv preprint (2026), version 2; accepted for RANDOM 2026. \href{https://doi.org/10.48550/arXiv.2602.09049}{doi:10.48550/arXiv.2602.09049}; \href{https://arxiv.org/abs/2602.09049}{arXiv:2602.09049}.

  \item[\textup{[CX26]}]\label{ref:gvu-random-regimes}
  T.-W. Chao and Z. Xu, ``Vertex-Minor Universality of a Random Graph,'' arXiv preprint (2026), version 3, 5 May 2026. \href{https://doi.org/10.48550/arXiv.2603.13600}{doi:10.48550/arXiv.2603.13600}; \href{https://arxiv.org/abs/2603.13600}{arXiv:2603.13600}.
\end{enumerate}

\subsection*{Comment}

No deterministic polynomial-time construction with $O(k^2)$ vertices was found in the public literature checked through 9 September 2026; the unresolved gap is between quadratic probabilistic existence and quartic efficient deterministic constructions. Exhaustive finite search can eventually find a quadratic-size graph, but that observation does not meet the polynomial-time requirement, and the resolved random-graph conjecture must not be relabelled as this open construction problem. The status audit used public primary sources and later-work searches; it is not an exhaustive citation-index audit. The related Pauli-pairability question asks only for perfect matchings on prescribed terminals, rather than all graphs, and does not require an efficient deterministic construction.

\subsection*{Field}

Quantum Resource Theory; Quantum algorithm

\subsection*{Topic}

Quantum state preparation; Resource conversion; Computational complexity and computability

\subsection*{ID}

\texttt{op\_4154bf4cbe0d288e}
